% ---------------------------------------------------------------------------
% PLOS Computational Biology submission
% From fluid shear to sclerostin: a mechanotransduction-resolved multiscale
% model of bone adaptation, and what it falsifies
%
% Built on the official PLOS LaTeX template (plos_latex_template.tex).
% Compiles standalone: the bibliography is a manual thebibliography, so
% plos2015.bst is NOT required.  A references.bib is shipped alongside for
% anyone who prefers BibTeX -- to use it, comment out the thebibliography
% block at the end and uncomment the two \bibliography lines.
%
% Figures are NOT embedded, per PLOS policy: captions sit in the text at
% first citation and the image files are uploaded separately (figures/Fig1
% ... Fig7).
% ---------------------------------------------------------------------------

\documentclass[10pt,letterpaper]{article}
\usepackage[top=0.85in,left=2.75in,footskip=0.75in]{geometry}

\usepackage{amsmath,amssymb}
\usepackage{changepage}
\usepackage[utf8]{inputenc}
\usepackage{textcomp,marvosym}
\usepackage{cite}
\usepackage{nameref,hyperref}
\usepackage[right]{lineno}
\usepackage{microtype}
\DisableLigatures[f]{encoding = *, family = * }

\usepackage[table]{xcolor}
\usepackage{array}

\newcolumntype{+}{!{\vrule width 2pt}}
\newlength\savedwidth
\newcommand\thickcline[1]{%
  \noalign{\global\savedwidth\arrayrulewidth\global\arrayrulewidth 2pt}%
  \cline{#1}%
  \noalign{\vskip\arrayrulewidth}%
  \noalign{\global\arrayrulewidth\savedwidth}%
}
\newcommand\thickhline{\noalign{\global\savedwidth\arrayrulewidth\global\arrayrulewidth 2pt}%
\hline
\noalign{\global\arrayrulewidth\savedwidth}}

% Uncomment for double spacing if the editor requests it
%\usepackage{setspace}
%\doublespacing

\raggedright
\setlength{\parindent}{0.5cm}
\textwidth 5.25in
\textheight 8.75in

\usepackage[aboveskip=1pt,labelfont=bf,labelsep=period,singlelinecheck=off]{caption}
\renewcommand{\figurename}{Fig}

% Remove brackets from numbering in List of References
\makeatletter
\renewcommand{\@biblabel}[1]{\quad#1.}
\makeatother

\usepackage{lastpage,fancyhdr,graphicx}
\usepackage{epstopdf}
\pagestyle{fancy}
\fancyhf{}
\rfoot{\thepage/\pageref{LastPage}}
\renewcommand{\headrulewidth}{0pt}
\renewcommand{\footrule}{\hrule height 2pt \vspace{2mm}}
\fancyheadoffset[L]{2.25in}
\fancyfootoffset[L]{2.25in}
\lfoot{\today}

% Convenience macro for microstrain
\newcommand{\ue}{\ensuremath{\mu\varepsilon}}

\begin{document}
\vspace*{0.2in}

\begin{flushleft}
{\Large
\textbf\newline{From fluid shear to sclerostin: a mechanotransduction-resolved multiscale model of bone adaptation, and what it falsifies}
}
\newline
\\
\textbf{Short title:} A mechanotransduction-resolved multiscale model of bone adaptation
\newline
\\
Hsiang-Lin Lee\textsuperscript{1,2},
Tzu-Ling Wang\textsuperscript{3},
Ming-Yu Hsieh\textsuperscript{4,5,6*}
\\
\bigskip
\textbf{1} Institute of Medicine, Chung Shan Medical University, Taichung, Taiwan
\\
\textbf{2} Department of Surgery, Chung Shan Medical University Hospital, Taichung, Taiwan
\\
\textbf{3} School of Nursing, Chung Shan Medical University, Taichung, Taiwan
\\
\textbf{4} School of Medicine, Chung Shan Medical University, Taichung, Taiwan
\\
\textbf{5} Division of Pediatric Surgery, Chung Shan Medical University Hospital, Taichung, Taiwan
\\
\textbf{6} Center for Evidence-Based Medicine, Chung Shan Medical University Hospital, Taichung, Taiwan
\\
\bigskip
* cshy1392@csh.org.tw
\end{flushleft}

% ---------------------------------------------------------------------------
\section*{Abstract}

Mechanical loading and calcium availability are the two levers most often
invoked in skeletal health, but no model has been able to weigh them against
each other, because mechanotransduction models terminate at the osteocyte and
bone-remodelling models contain no mechanics. We built an eight-module
multiscale model that carries a load continuously from organ-level bending,
through poroelastic lacunocanalicular shear, through three-state
mechanosensitive channel gating, through
Piezo1--Ca\textsuperscript{2+}--YAP/TAZ--sclerostin--Wnt signalling, through
the RANK/RANKL/OPG cell populations, to three-surface cortical geometry,
mineralisation and areal bone mineral density, with systemic
calcium--parathyroid hormone--vitamin D coupled bidirectionally. Loading
enters strictly as force; strain is a regulated output, preserving the
mechanostat feedback loop. Six free parameters were fitted to six targets;
eight further observations were held out and never entered the objective. The
model reproduces adult turnover (7.33\,\%/yr), disuse loss (1.15\,\%/month),
the non-progressive plateau of calcium supplementation, and romosozumab's
spine response ($+12.5\,\%$ at 12 months) and post-withdrawal loss
($-12.4\,\%$). Frost's mechanostat set point emerges at 787\,\ue, inside the
lazy zone bracketed by the remodelling (100--300\,\ue) and modelling
(1500--3000\,\ue) thresholds he inferred from observation, without being
fitted. Loading raises bone mineral content 15.1-fold more than calcium, but
only at a specified intensity ($\approx$2900\,\ue) and in bone mass rather
than areal density, which dilutes the response by 61\,\%. Unilateral loading
reproduces all six peripheral quantitative computed tomography indices of the
tennis-player humerus, while systemic interventions produce side-to-side
differences of 0.0000\,\%. Our starting hypothesis that osteoporosis is
an alternative stable state is falsified: the model is monostable, and a full
withdrawal-and-recovery probe returns to 100.3\,\% of baseline with no
hysteresis. Skeletal irreversibility must therefore arise from structural loss
outside continuum porosity dynamics. One of the three clauses of our
loading-versus-calcium prediction reverses sign under a different densitometric
measure, which we report as a methodological result in its own right.

% ---------------------------------------------------------------------------
\section*{Author summary}

Bone adapts to the forces placed on it. Two ingredients are usually named:
calcium, the material bone is built from, and mechanical loading, the signal
that tells bone where to build. Clinicians and patients often treat these as
interchangeable levers, but they are not the same kind of thing, and no
computational model has been able to put them in the same framework and
compare them --- because the models that describe how a bone cell senses force
stop before bone density, and the models that predict bone density have no
real force in them.

We joined the two ends. Our model follows a load all the way from the bending
of a whole bone, through fluid squeezed past the cells buried inside it,
through the ion channels those cells open, through the molecular signals they
release, to the cell populations that add and remove bone, and finally to the
density a clinical scan would measure. Because the chain is unbroken, we could
ask questions that neither half could answer alone.

Three results matter. Loading raises bone mass about ten times more than
calcium does, but only when both the loading intensity and the measurement are
stated precisely --- the same model gives opposite answers under different
measures, which is a warning for how such claims are reported. Purely local
loading produces the side-to-side asymmetry seen in tennis players, and no
whole-body treatment can imitate it. And our own starting hypothesis --- that
osteoporosis is a second stable state a skeleton can fall into and cannot
climb out of --- is wrong in this framework: the model recovers fully.
Irreversibility must come from structural damage the model does not yet
represent.

\linenumbers

% ---------------------------------------------------------------------------
\section*{Introduction}

Bone is a structure that carries its own sensors. Frost's mechanostat
\cite{frost1987} proposed that bone tissue holds strain near a set point by
adding material when strain is high and removing it when strain is low, and
the framework has organised four decades of skeletal physiology. What it has
never had is a mechanism: the set point is inferred from observation, not
derived, and the theory is silent on the molecular chain that would produce
it.

That chain is now largely known. Load deforms the bone matrix, driving
interstitial fluid through the lacunocanalicular network; the resulting shear
stress on osteocyte processes gates mechanosensitive ion channels including
Piezo1 \cite{li2019}; the calcium transient that follows shifts YAP/TAZ to the
nucleus; sclerostin, the osteocyte-secreted Wnt antagonist encoded by
\emph{SOST}, falls; Wnt/$\beta$-catenin signalling rises; and osteoblast and
osteoclast populations respond. Each link is supported, and several are drug
targets --- romosozumab, an anti-sclerostin antibody, is licensed on exactly
this axis \cite{cosman2016}.

Computational work has advanced from both ends and met in the middle nowhere.
On the sensing side, Fu and colleagues \cite{fu2025} built a whole
bone--lacunocanalicular--osteocyte model that reproduces the effects of
loading magnitude, frequency, cycle number and recovery time, but stops at the
osteocyte response. On the tissue side, the Lemaire \cite{lemaire2004} and
Pivonka \cite{pivonka2008} families of models describe RANK/RANKL/OPG
population dynamics and drug interventions in quantitative detail, but drive
them with a phenomenological stimulus rather than mechanics; where mechanics
appears at all it is as strain energy density imposed from outside
\cite{scheiner2013}. The consequence is that no existing model can (i)
distinguish two loading protocols with identical energy but different waveform
and report a bone mineral density trajectory, (ii) compare mechanical loading
against an anti-sclerostin antibody in the same currency, or (iii) let calcium
intake and mechanical loading compete for the same bone.

The third gap is the clinically loudest. Calcium supplementation produces a
small, non-progressive increment in bone density --- Tai and colleagues'
meta-analysis \cite{tai2015} found $+0.7$ to $+1.8\,\%$, reached within one
year and not continuing thereafter --- while exercise interventions and, most
starkly, the natural experiment of unilateral racquet sports
\cite{haapasalo2000} produce far larger and site-confined changes. The natural
reading is that calcium is \emph{permissive} and loading is \emph{instructive}:
calcium determines whether bone can be built, loading determines whether it
will be. That reading has never been tested inside a single mechanistic model,
because no model has contained both.

We report such a model, and three falsifiable predictions. All three were
written into the project plan before any parameter was fitted, and the dated
plan and its decision log are in the archived repository, so the ordering can
be checked; we make no claim of formal pre-registration in a public registry. \textbf{P1}: in the calcium-replete range, the marginal effect
of calcium on bone density is small, while appropriate loading is large, and
the two interact non-additively. \textbf{P2}: a two-compartment model sharing
one systemic calcium pool reproduces the tennis-player asymmetry, and no
systemic intervention can. \textbf{P3}: with mechanosensitivity, sclerostin
secretion or oestrogen as bifurcation parameters, the system exhibits a
saddle-node bifurcation and bistability, so that osteoporosis is an
alternative stable state --- which would explain why bone is lost quickly and
regained slowly.

Two of these survive. The third does not, and we report its falsification in
the main text because it is the strongest available evidence that the model
predicts rather than fits.

% ---------------------------------------------------------------------------
\section*{Results}

\subsection*{An unbroken chain from force to bone mineral density}

The model comprises eight modules (M1--M8; \nameref{S1_Text}). M1 computes
cross-sectional strain from a prescribed peak bending moment and axial force.
M2 solves Biot poroelasticity for lacunocanalicular fluid shear
\cite{weinbaum1994}. M3 integrates three-state (open/closed/inactivated)
mechanosensitive channel gating \cite{fu2025} over a day's loading to produce
a single daily scalar dose. M4 and M5 carry
Piezo1--Ca\textsuperscript{2+}--YAP/TAZ--sclerostin--Wnt signalling. M6 is a
RANK/RANKL/OPG osteoblast--osteoclast population model
\cite{lemaire2004,pivonka2008}. M7 evolves periosteal, endocortical and
intracortical surfaces plus mean tissue mineral density. M8 is systemic
calcium--parathyroid hormone--1,25(OH)\textsubscript{2}D
\cite{peterson2010}, coupled bidirectionally. The single-compartment state
vector has 17 elements; the two-compartment version has 30.

One interface decision determines whether the model is a mechanostat at all.
\textbf{Loading enters only as force} --- peak bending moment and axial force
--- and strain is computed as an output of current geometry and material
state. Prescribing strain, as an earlier version of this model did, silently
severs the negative feedback: bone volume fraction becomes a pure integrator,
the calcium plateau cannot appear, and there is no set point for bifurcation
analysis to find. We verified the closed loop numerically:
$\mathrm{d}\ln\varepsilon_p/\mathrm{d}\ln f_{bm} = -3.1300$ in cortex and
$-2.0000$ in the vertebral compartment, each matching its own $-\kappa_E$ to
better than $4\times10^{-16}$, so added bone does reduce strain. The two
exponents differ because compact and trabecular bone are different materials;
we return to this in \nameref{S1_Text}.

Two independent cross-validations support the mechanics. A Crank--Nicolson
finite-difference solution of the poroelastic problem agrees with the
closed-form Biot solution to a maximum relative error of 0.0029\,\%, and an
affine per-cycle channel operator agrees with stepwise integration to
$2.2\times10^{-15}$.

\subsection*{Calibration, and what was withheld}

Five parameters were fitted against five targets: adult turnover, disuse loss,
calcium effect size, calcium plateau slope, the romosozumab spine response,
and the post-withdrawal resorption-marker overshoot. The formation rate
constant is not free --- it is derived at every parameter set from bone-mass
balance, since the three-surface geometry admits no stationary geometric fixed
point. All other parameters come from the literature (\nameref{S1_Table}).

Eight observations were \textbf{held out} and never entered the objective: the
six peripheral quantitative computed tomography indices of the tennis-player
humerus, the post-withdrawal bone density loss, and the emergent mechanostat
set point. Their status as blind tests is what the rest of this section rests
on.

Identifiability analysis (\nameref{S1_Fig}) shows \texttt{k\_res} and
\texttt{K\_S} are individually identifiable, each pinned by one target with a
clear single minimum. The antibody clearance constant \texttt{delta\_ab} is
\textbf{not}: its conditional $\chi^2$ falls monotonically to its upper bound,
because the objective's midpoint metric pulls the romosozumab response toward the
centre of its band and the parameter cannot reach that far. The value we report
satisfies the band; it is not a point estimate, and we do not treat it as one.
The two PTH-pathway parameters are \textbf{practically non-identifiable}:
76.9\,\% of the scanned grid lies inside the confidence region and the ridge
correlation is only $-0.26$, so they are not trading off along a ridge --- both
are simply weakly constrained. We report the joint region rather than marginal
intervals, and we do not interpret either parameter's point value.
\texttt{K\_P\_sost} no longer sits at its upper bound, as it did before the
osteocyte turnover constant was corrected, but it remains high, and the
interpretation is unchanged: large values mean PTH inhibition of \emph{SOST} is
weak and the fit is asking for that pathway to be nearly absent, as Lemaire and
Pivonka assume by omitting it.

\subsection*{The mechanostat set point is an output, and it lands where Frost put it}

At baseline the model sits at a stable fixed point
($\max\mathrm{Re}\,\lambda = -2.4\times10^{-4}$), with 24-month areal bone
mineral density drift of $-0.015\,\%$ and turnover of 7.33\,\%/yr
(Fig~\ref{fig1}). The bar was 5--10\,\%/yr, the upper part of the
3--10\,\%/yr range conventionally quoted for cortical bone
\cite{parfitt2002}; our value sits in the upper half of it, so this is a pass
against a target we set at the fast end of what is quoted.

\begin{figure}[!h]
\caption{{\bf Baseline steady state and the emergent mechanostat set point.}
(A) Areal bone mineral density over 24 months at baseline; drift
$-0.015\,\%$. (B) Bone volume fraction over the same period. (C) The peak
strain the closed loop settles at, 786.8\,\ue, against the 300--1500\,\ue
band Frost inferred from observation. This value is computed, not fitted, and was held
out of calibration.}
\label{fig1}
\end{figure}

The peak strain the closed loop settles at is \textbf{786.8\,\ue}. Frost
\cite{frost1987} inferred two separate thresholds from observation: strains at
or above 1500--3000\,\ue{} drive modelling that adds cortical bone, and strains
below 100--300\,\ue{} release remodelling that removes it. The emergent value
sits inside the lazy zone those thresholds bracket --- above the remodelling
threshold, below the modelling threshold --- and it is computed here from
channel gating kinetics and signalling Hill coefficients, without ever being
fitted. This is the model's most direct evidence that a phenomenological law
can be reduced to molecular mechanism, and it costs no free parameter.

The same run also shows why cortical bone gains size rather than density: at a
bone volume fraction of about 0.95 the cortex is within 5\,\% of its
densification ceiling, and
$\mathrm{d}\ln\varepsilon/\mathrm{d}\ln r_p = -4.27$ is substantially stronger
than $\mathrm{d}\ln\varepsilon/\mathrm{d}\ln f_{bm} = -3.13$. Growing outward
is simply the cheaper way to reduce strain --- an emergent consequence of
parameter values, not an imposed rule.

\subsection*{Loading waveform: two targets recovered, one sign disagreement}

Fig~\ref{fig2} shows three orthogonal slices through loading-waveform space.

\begin{figure}[!h]
\caption{{\bf Mechanical dose--response.}
(A) Cycle number against 12-month areal density change; marginal return falls
33-fold from 18 to 1200 cycles. The saturation is a downstream property ---
at the channel level the dose is linear in cycle number. (B) Frequency at
both levels: fluid shear (grey, left axis) rises with $\ln f$ as
poroelasticity requires, while the bone-level response (teal, right axis)
falls. The logarithmic form survives to bone; the sign does not. (C) Rest
insertion at three amplitudes. The gain falls with amplitude, matching
experiment, and reverses the direction seen at the channel level.}
\label{fig2}
\end{figure}

\emph{Cycle number.} Increasing cycles from 18 to 1200 raises the 12-month
response from $+0.503\,\%$ to $+2.286\,\%$, with marginal return falling
33-fold and monotonically. This reproduces the saturation of the osteogenic
response beyond a few dozen cycles \cite{yang2017}. Notably, the saturation is
\textbf{not} in the channel: at the channel level the dose is asymptotically
linear in cycle number across four orders of magnitude of gating rate
constants, because periodic driving puts the three-state system on a limit
cycle where every subsequent cycle contributes identically. The saturation is
a downstream property of the sclerostin--Wnt--cell cascade, and can only be
seen at the bone level.

\emph{Rest insertion.} The target here is conditional: a substantial gain at
low amplitude that fades as amplitude rises. Inserting 10\,s between cycles
amplifies the response 4.60-fold at low amplitude, 1.66-fold at moderate and
1.45-fold at high. The \textbf{direction} matches Srinivasan and
colleagues \cite{srinivasan2002} and the \textbf{magnitude is now short of it}:
we obtain 4.60-fold against their measured 5.8-fold --- labelled periosteal
surface of 21.9\,\% with rest insertion against 3.8\,\% without, at a matched
820--830\,$\mu\varepsilon$ --- a shortfall of 21\,\%. We report the history
because it is a caution against reading close agreement as validation: this
quantity read 13.99-fold before the elastic-law exponents were taken from
Currey \cite{currey1988}, 6.02-fold after that, and 4.60-fold once the
osteocyte turnover constant was corrected (see Limitations). Three successive
corrections, none of them aimed at this target, moved it past the measured
value and out the other side. The conditional \emph{shape} the target is really
about --- a large gain at low amplitude that fades as amplitude rises --- has
survived all three. We
flag that this agreement is recent and was not engineered: until the elastic-law
exponents were taken from Currey the same quantity was 13.99-fold, and the band
we had been scoring against, 2--5-fold, was one we invented rather than read off
the source. This too reverses between
levels: at the channel level our gain rose with amplitude, contradicting
experiment. Downstream saturation flips it to the observed direction, which is
a non-trivial success for the cascade structure.

\emph{Frequency.} Here the model disagrees with experiment, and we state it
plainly. Fluid shear rises with frequency exactly as poroelasticity requires
($\tau$ 4.62 to 20.02\,Pa over 1--10\,Hz; $r$ against $\ln f = +0.999$, with
no fitted parameter). The logarithmic \textbf{form} survives to the bone level
--- consistent with time-lapsed micro-computed tomography estimates of
mechanostat parameters \cite{marques2023} --- but the \textbf{sign does not}:
areal density change falls from $+0.715\,\%$ to $+0.560\,\%$ ($r = -0.997$). A
duration-matched control, holding bout length fixed while cycle number scales
with frequency, gives the same sign, so this is not a session-length artefact.
The mechanism is that the channel opening rate saturates in shear, so the
shorter loaded interval per cycle at high frequency costs more than the higher
peak shear gains. Experimentally, higher frequency within this range is more
osteogenic \cite{hsieh2001}. This is the sharpest falsifiable consequence of
our placeholder channel rate constants and the first thing to re-examine when
better values are available.

\subsection*{P1: calcium is permissive, loading is instructive --- under stated measures}

Fig~\ref{fig3} shows a factorial of four calcium intakes against three loading
levels over 24 months.

\begin{figure}[!h]
\caption{{\bf Calcium against loading (P1).}
(A) 24-month bone mineral content change across four calcium intakes and three
loading levels; dashed line marks P1's 4\,\% threshold. Only the vigorous
programme (2949\,\ue) clears it. (B) The same vigorous-loading contrast in
bone mineral content and in areal density: dual-energy X-ray absorptiometry
dilutes the loading response by 61\,\%, because loading moves the periosteal
surface outward and areal density divides by projected width.}
\label{fig3}
\end{figure}

P1 has three clauses, and we state the bar for each before giving the numbers.
\textbf{Clause 1}: within the calcium-replete range, the marginal effect of
calcium on 24-month bone density is under 1\,\%. \textbf{Clause 2}: an
appropriate loading programme exceeds 4\,\% over the same period.
\textbf{Clause 3}: the two interact non-additively, calcium deficiency
truncating what loading can build.

\emph{Clause 1.} Raising calcium from 800 to 1500\,mg/day --- the replete
range the clause is about --- changes 24-month areal density by
\textbf{$+0.304\,\%$}, under the 1\,\% bar. Across the wider 400 to
1500\,mg/day contrast against which calibration target V7 is defined, the
change is \textbf{$+0.663\,\%$, just below the $+0.7$ to $+1.8\,\%$ band}
Tai and colleagues report \cite{tai2015} --- a miss of 0.037 against a declared
tolerance of 0.8, and one we report rather than fit away. It arrived when the
osteocyte turnover constant was corrected (see Limitations), and its direction is
the informative part: the calcium arm had been riding on an artefact, and
removing the artefact made calcium \emph{less} effective. In both contrasts the increment reaches a
plateau within the first year rather than accumulating. That non-progressive
\emph{shape} is the discriminating feature rather than the magnitude, and it is
a fitted target in its own right with its own bar: the slope of the increment
over the second half of the run must be under 0.3\,\%/yr in absolute value,
and it is $-0.131\,\%$/yr --- flat, and if anything very slightly
declining.

\emph{Clause 2.} An appropriate loading programme --- which we define
explicitly as peak strain $\approx$2900\,\ue, still below the 3000\,\ue{}
physiological ceiling and comparable to an elite tennis player's humerus ---
raises bone mineral content by \textbf{$+4.582\,\%$}, over the 4\,\% bar.
The ratio to calcium is \textbf{15.1-fold}.

Two qualifications are not caveats but results.

\textbf{The measure changes the answer.} The same loading contrast measured as
areal density is $+1.778\,\%$, not $+4.582\,\%$: areal density divides bone
mass by projected width, and loading moves the periosteal surface outward, so
\textbf{dual-energy X-ray absorptiometry dilutes the loading response by
61\,\%}. Our model retains this size artefact deliberately, because it is what
allows one model to be read against both the densitometry literature and the
peripheral quantitative computed tomography literature. Reporting clause P1-2
in areal density would report a measurement artefact as biology.

\textbf{The intensity changes the answer.} Moderate resistance loading
(2236\,\ue) gives only $+3.118\,\%$ and does not support the claim.
``Appropriate loading'', left undefined, makes the prediction unfalsifiable,
since any shortfall can be attributed to insufficiently appropriate loading.

\emph{Clause 3} is supported on both framings, and we flag that this was not
true of an earlier version of the model. Calcium deficiency truncates what
loading can build: under vigorous loading, severe deficiency achieves
$+3.994$\,\% against $+4.756$\,\% when supplemented, so deficiency costs
16\,\% of the gain. The \emph{difference-in-differences} interaction --- loaded
minus sedentary, compared across calcium levels --- is \textbf{$+0.085$}
percentage points, the same sign, so loading's marginal benefit really is larger
when calcium is adequate.

We report the history because it bears on how much weight this clause can carry.
Before the systemic calcium module was rebuilt (see Limitations), the same
quantity was \textbf{$-0.536$} percentage points --- the opposite sign ---
because deficiency dragged the sedentary comparator down faster than the loaded
arm, a floor effect in the comparator rather than synergy. The correction that
removed a supraphysiological swing in serum calcium also flipped this
interaction. A conclusion that changes sign when an unrelated module is repaired
is one to hold loosely.

\textbf{One of P1's three clauses therefore reverses sign under a different but
equally reasonable measure}, and a second did so until a module elsewhere in the
model was corrected. We regard this as a substantive finding about how
such claims must be reported, not a weakness of the model, and we return to it
in the Discussion.

\subsection*{P2: only local loading can produce a local difference}

Fig~\ref{fig6} shows the two-compartment result: one humerus loaded, one not,
sharing a single systemic calcium and parathyroid hormone pool, over five
years.

\begin{figure}[!h]
\caption{{\bf Site specificity (P2).}
(A) The six peripheral quantitative computed tomography indices of the loaded
humerus against Haapasalo's measured ranges (grey bands). All six fall within
tolerance, and all six were held out of calibration. (B) Side-to-side
difference by intervention type: only unilateral loading produces one;
systemic calcium and systemic romosozumab produce 0.0000\,\%.}
\label{fig6}
\end{figure}

P2 has two halves, and the second is the one that can fail cleanly.
\textbf{First}, unilateral loading should reproduce the racquet-sport
humerus. \textbf{Second}, no systemic intervention should produce a
side-to-side difference at all --- a single non-zero asymmetry under a
systemic agent would falsify it.

All six indices reported by Haapasalo and colleagues \cite{haapasalo2000} fall
within tolerance, model against measured range: bone mineral content
$+12.38\,\%$ against 14--27\,\%, total area $+12.17\,\%$ against
16--21\,\%, cortical area $+10.40\,\%$ against 12--32\,\%, maximal second
moment of area $+24.48\,\%$ against 27--67\,\%, medullary cavity
$+13.93\,\%$ against 19\,\%, and cortical volumetric density $+1.79\,\%$
against $-2$ to $0\,\%$. Per-index tolerances are listed in
\nameref{S2_Table}. These are \textbf{hold-out} values. The discriminating one is the last: the
playing arm's gain is geometric, with volumetric density essentially
unchanged, and the model was never fitted to that constraint.

The complementary claim is the sharper one. With both arms loaded identically
and a systemic intervention applied --- 1500\,mg/day calcium, or romosozumab
--- the side-to-side difference is \textbf{$+12.38$\,\%} for loading against
\textbf{$-0.0000$\,\%} and \textbf{$-0.0000$\,\%} respectively. Because both compartments see the
identical shared parathyroid hormone, no systemic agent can manufacture a
local difference. The tennis-player asymmetry is not reproducible by any pill,
in this framework as a matter of structure rather than of parameter values.

\subsection*{Disuse and the minimum effective countermeasure}

Unloading produces 1.220\,\%/month bone loss over the 180-day window
(Fig~\ref{fig4}). The bar was the 1.0--1.5\,\%/month band, which sits inside
the 1--2\,\%/month decline reported at weight-bearing sites over 90 days of
head-down bed rest \cite{spatz2012}. The quantity we compute is bone mineral
content rather than areal density; in disuse the periosteum barely moves
($-0.002$\,mm over 180 days) and the two agree to 0.003 percentage points per
month, unlike the loading contrast above where they differ by 61\,\%.

Sclerostin rises, which is the direction the target requires, but the
\emph{shape} is wrong and we report this as a partial match. Serum sclerostin
in the model reaches $+58.0\,\%$ and then stays there, essentially flat from
day 28 onward. Measured in men over 90 days of bed rest it rises to
$+29\,\%$ at day 28, peaks at $+42\,\%$ at day 60, and has fallen back to
$+22\,\%$ by day 90 \cite{spatz2012}. We reproduce the rise and roughly its
peak magnitude, but not the attenuation: sclerostin here equilibrates within
days to the new mechanical dose and has no adaptive decline.

The same experiment supplies a second check, on the systemic arm, and it is one
the model failed until we repaired it. In those subjects parathyroid hormone
fell 17\,\% by day 28 and 24\,\% by day 60, with urinary calcium
significantly elevated throughout \cite{spatz2012} --- the classic
immobilisation picture, in which calcium released from unloaded bone suppresses
the parathyroid. The model now reproduces it: parathyroid hormone falls
\textbf{21.2\,\%} by day 28 and \textbf{21.3\,\%} by day 60, on a serum
calcium rise of 2.8\,\%.

It did not before. In the version of this module we first built, the same run
moved parathyroid hormone by 0.01\,\% --- the coefficients carrying skeletal
calcium into the serum pool were three orders of magnitude below the flux they
were normalised against, so bone could not move serum calcium at all, and the
coupling we describe as bidirectional was in practice one-way. We report this
because the failure was invisible to every target we had: nothing in the
validation set asked the model to move systemic calcium from bone, so nothing
caught it until we tested a fact the literature had recorded and we had not.

\begin{figure}[!h]
\caption{{\bf Disuse and the minimum effective countermeasure.}
(A) Bone mineral content over 180 days of unloading, unprotected and with the
best in-budget countermeasure. (B) Monthly loss across the countermeasure
grid; blank cells exceed the 30\,min/day budget. No in-budget regimen holds
loss below 0.25\,\%/month.}
\label{fig4}
\end{figure}

Searching loading prescriptions under a $\le$30\,min/day constraint, the best
regimen within budget (peak moment $\times$4.5, 80 cycles, 0.5\,Hz, 10\,s
inter-cycle rest, 16.0\,min/day) reduces loss to 0.771\,\%/month ---
preventing 37\,\% of it. \textbf{No regimen within the time budget holds loss
below 0.25\,\%/month.} We report this negative result because it is clinically
relevant: in this model, time-efficient loading alone does not arrest disuse
osteopenia.

We also flag a scope limit that our validity check surfaced. Beyond about 7.5
months, complete unloading drives the model to its porosity floor and out of
the linear-elastic domain the entire mechanics stack assumes. Real bed rest
decelerates and plateaus at 30--50\,\% loss. All disuse results here are
confined to the 180-day window on which the target is defined, and every run
reports whether it stayed inside the elastic domain.

\subsection*{Pharmacology: self-limitation, withdrawal rebound, and a shared node}

Fig~\ref{fig5} shows four arms over 24 months with romosozumab stopped at
month 12.

\begin{figure}[!h]
\caption{{\bf Romosozumab, withdrawal and loading.}
(A) Four arms over 24 months, drug stopped at month 12. (B) The osteoblast
population on drug: it peaks early and declines while dosing continues,
because compensatory \emph{SOST} transcription progressively offsets the
antibody. (C) Trabecular against cortical response, with the $+11$--$14\,\%$
spine target band shaded; the 5.3-fold amplification comes from structure
alone.}
\label{fig5}
\end{figure}

Because the $+11$--$14\,\%$ romosozumab target is a \emph{lumbar spine} number
\cite{cosman2016} and the spine is trabecular, we evaluate it in a trabecular
vertebral compartment derived from the cortical parameter set by changing only
structure, geometry and load --- the biology is shared. One consistency check does
not come free, and saying why is more useful than the check. Until v2.16 a
single exponent served both compartments; taking the measured value shows that
agreement was a compromise, not a validation. Currey's power law over compact
bone \cite{currey1988} gives 3.13, which at a bone volume fraction of 0.12 puts
the vertebra at 26\,MPa --- below the measured 50--300\,MPa range. Trabecular
bone is a cellular solid, and Gibson--Ashby open-cell scaling gives 2. Using
3.13 in cortex and 2.00 here gives an apparent modulus of 288.0\,MPa, inside
the measured range, on a measurement at each end rather than one value split
between them.

In that compartment romosozumab gives \textbf{$+12.459\,\%$} at 12 months,
inside the 11--14\,\% target band, against $+2.334\,\%$ in cortex --- a
5.3-fold amplification arising from structure alone. This is a
\textbf{calibration} target, not a prediction: \texttt{delta\_ab} was fitted to
it. The amplification factor, by contrast, has no direct empirical counterpart
and we do not present it as validated. The nearest clinical comparison is the
spine-to-hip contrast, $+15.1\,\%$ against $+5.4\,\%$ at 24 months on the
210\,mg monthly regimen \cite{mcclung2018}, about threefold --- but the total
hip is a mixed cortical--trabecular site whereas our cortical compartment is a
pure diaphysis, so a larger ratio is expected and the two numbers are not like
for like. Building the compartment also
forced a general constraint into the open: because the signalling chain is
shared, both compartments must sit at the same emergent shear set point, which
a vertebra at realistic load and bone volume fraction does not do at the same
tissue strain. The poroelastic transfer coefficient must therefore differ
between trabecular and osteonal bone --- defensible, since it is a
microstructural property, and set here by matching baseline shear rather than
by fitting any outcome.

\emph{Self-limitation.} The formation marker rises, peaks, and then declines
while dosing continues (rise $+0.421$ in the first half of the course against
$-0.259$ in the second), which is the shape the marker data show. The
\textbf{timing does not match}: our marker peaks at about month 0.2 and
reaches 3.169-fold baseline, whereas formation markers clinically peak near
month 1 and decline towards baseline over 6--12 months. We reproduce the
self-limitation but compress its onset. The mechanism is explicit: compensatory
up-regulation of \emph{SOST} transcription accumulates under antibody exposure
and progressively offsets it, so free sclerostin climbs back toward baseline
by month 9 while treatment is still running.

\emph{Withdrawal.} That same compensation decays on its own clock rather than
the antibody's, so when clearance stops the raised production does not,
sclerostin overshoots, and resorption surges through the RANKL-dependent arm
of sclerostin action \cite{wijenayaka2011}. We calibrated this mechanism
against the post-withdrawal \textbf{resorption marker} overshoot (1.124
$\times$ baseline, \textbf{just outside} the 1.2--1.4 range corresponding to the rise of
$\beta$-CTX above baseline after transition to placebo \cite{mcclung2018}),
deliberately not against bone density --- which leaves the density fall a blind
prediction.

The bar we set for that prediction was \textbf{direction only}: bone density
must fall over the washout. It does, by \textbf{$-8.70\,\%$} over the 12
months after the last dose, and we state plainly that a directional bar is a
weak one. Judged on magnitude rather than sign, this hold-out is now the closest
it has been. The model gives back \textbf{78.5\,\%} of the treatment gain and
finishes \textbf{2.68\,\% above} pre-treatment baseline at 24 months, against
the observation that bone density returned toward pretreatment levels on placebo
\cite{mcclung2018} --- toward baseline, not through it, which is what the model
now does. We note that this improved when the osteocyte turnover constant was
corrected and no drug parameter was touched: before that correction the same run
overshot, finishing 0.75\,\% \emph{below} baseline. A hold-out moving toward
the data under a change made for an unrelated reason is the kind of evidence we
weight most heavily, and we flag equally that the remaining discrepancy runs the
other way --- real washout is slower than a step function, because the antibody
has a pharmacokinetic profile the model does not.

The resorption-marker target itself is now \textbf{missed}, at 1.124 against a
band of 1.2--1.4 --- but we must qualify what that means, because the band is
ours rather than the literature's. McClung and colleagues report only that
$\beta$-CTX increased above baseline after transition to placebo
\cite{mcclung2018}, a direction with no magnitude, and the FRAME trial cannot
supply one because its patients transitioned to denosumab rather than placebo
\cite{cosman2016}. The sourced content of this target is therefore ratio $> 1$,
which the model meets. What follows concerns the band we set ourselves, and the
reason we still report the shortfall is that the underlying conflict is real
regardless of where the band sits.
Rebuilding the systemic calcium module (see Limitations) made skeletal calcium
flux large enough to move the parathyroid, and that creates a conflict: the same
PTH$\rightarrow$RANKL gain that lets dietary calcium protect bone --- which
target V7 requires to be strong --- also lets the calcium flood released after
withdrawal suppress PTH and damp the very resorption overshoot V16 measures.
Across the admissible range of that gain, V7 needs it high and V16 needs it low,
and no value satisfies both. We chose the midpoint, and both targets now sit
outside their bands --- V7 short by 0.037 against a tolerance of 0.8, V16 short
by 0.076 against a tolerance of 0.15. That both fall on the same side is worth
stating: the conflict is not a trade we won on one side, it is a constraint
neither side satisfies. The
decoupled module we used previously could not have exposed this conflict,
because in it bone could not move serum calcium at all.

\emph{Interaction.} Loading combined with the antibody gives $+2.826\,\%$ at
12 months against $+2.334\,\%$ and $+0.765\,\%$ alone: positive and larger
than either, but \textbf{8.8\,\% sub-additive}, because both act through the
same sclerostin--Wnt node and partly compete for it. Schulte and colleagues
\cite{schulte2026} report loading to be additive and synergistic with anabolic
agents and not with antiresorptives. Our qualitative direction agrees for the
anabolic arm; our strict additivity does not. \textbf{The antiresorptive half
of that contrast is untested here}: the model carries no bisphosphonate
mechanism, so we make no claim about it.

\subsection*{Four targets we had specified but never run}

Preparing this report we found that four of the twenty targets registered in
\textbf{S2 Table} --- the Mendelian-randomisation constraint on serum calcium
(V7b), the loss--recovery asymmetry (V11), the dual action of sclerostin on
resorption (V13), and the geometry of the postmenopausal cortex (V15) --- had
mechanisms built for them, and commented in the source as existing for them, but
had never been evaluated by any experiment. We ran them. One passes and three do
not, and we report all four, because a validation table listing targets nobody
checked is worse than a shorter table.

\emph{The loss--recovery asymmetry (V11) passes, and comfortably.} After 180
days of unloading, half the lost bone mineral content accrues in 140 days but
takes 2014 days to win back once loading resumes --- a ratio of 14.4 against a
target of 3 --- and ten years of reambulation recovers all of it and then some, ending 146.5\,\% of the loss recovered. This
is a \textbf{rate} asymmetry and not bistability: recovery is monotone and still
proceeding, which is consistent with, not in tension with, the falsification of
P3 below.

\emph{The Mendelian-randomisation constraint (V7b) fails on sign.} Clamping
serum calcium 2\,\% above baseline at fixed intake raises 24-month areal density
by $+0.668$\,\%, where the human genetic evidence requires no increase. The
magnitude is an order of magnitude below the loading response, so P1's
permissive-versus-instructive contrast survives; the sign does not. The same run
exposed a defect no target had been written to catch, and it was serious enough
that we rebuilt the module rather than report around it. Across 400 to
1500\,mg/day of intake, serum ionised calcium in the model moved by 15\,\%,
where physiology defends it within about 2\,\%: the transcellular absorption arm
had been written as a fraction rather than a flux and so contributed under
0.1\,\% of absorption, leaving the unregulated paracellular route to dominate,
and renal excretion was a bare power law rather than the saturable tubular
reabsorption it had been specified as. Our regression suite missed it because the
test guarding it was written as a runaway check with a 15\,\% bar and was
thereafter read as a homeostasis check. After rebuilding both arms the excursion
is \textbf{1.6\,\%}, and the V7b sign failure above is measured on the repaired
module. The episode is why we now report V7's effect size as sitting at the low
edge of the meta-analytic band \cite{tai2015} rather than comfortably inside it:
the earlier, larger value was riding on a serum-calcium swing that cannot occur.

\emph{The dual action of sclerostin (V13) is right in direction and short in
magnitude.} Over a twelve-fold sclerostin sweep, RANKL rises (1.00 to 1.29) and
OPG falls (1.00 to 0.44) as required, but osteoclast activation saturates at
and the RANKL:OPG ratio therefore rises \textbf{2.9-fold}, which is the quantity
the source reports rising dose-dependently \cite{wijenayaka2011}.

We had originally scored this target against the roughly 7-fold figure in the
same paper and recorded it as a failure. Reading the full text shows that was our
error, not the model's: the 7-fold is osteoclastic \textbf{resorption} --- pit
area in a 14-day osteocyte--splenocyte co-culture at 100\,ng/ml recombinant
sclerostin --- and not the instantaneous activation term we were comparing it
against. It is the same wrong-contrast mistake this section exists to expose,
committed by us. The activation term does carry a hard ceiling of 1.42-fold,
because its Hill constant sits below baseline sclerostin, but that ceiling
answers a question the source never asked. Scoring V13 properly means running the
co-culture's downstream resorption output, which we have not done.

\emph{The postmenopausal geometry (V15) fails on its second half, and the way it
fails is instructive.} At severely reduced oestrogen the endocortical radius
expands ($+0.0752$\,mm) and the cortex thins ($-0.0773$\,mm) as required, but the
periosteal radius drifts very slightly \textbf{inward} ($-0.0021$\,mm) rather
than outward, so the model does not reproduce the widening that accompanies
postmenopausal thinning. A ten-year run appears to pass this target
spectacularly, with 4.62\,mm of periosteal expansion --- but it leaves the
linear-elastic domain on day 511, and is the same class of artefact that the
modelling saturation bound was introduced to remove. Read inside the valid
domain, the target fails.

\subsection*{P3 is falsified: the model is monostable}

P3 held that osteoporosis is an \textbf{alternative stable state}: that past
some critical loss of oestrogen the skeleton settles into a second,
self-sustaining low-bone-mass equilibrium, and that removing the original
insult does not bring it back. That is what a saddle-node bifurcation
produces, and it is an attractive account of the clinical asymmetry between
rapid loss and slow, incomplete regain. Two positive feedbacks were built
into the model specifically because they could generate it: remodelling
surface is lost as bone disappears, and sensing capacity is lost as
osteocytes die.

Confirmation would have required three things: two stable equilibria
coexisting at the same oestrogen level, a fold at which one of them
vanishes, and hysteresis --- a path down that the system cannot retrace on
the way back up. We find none of the three (Fig~\ref{fig7}).

\begin{figure}[!h]
\caption{{\bf Osteoporosis is a steep continuous shift, not a second stable
state (P3).}
(A) The oestrogen continuation with envelope geometry frozen: a single branch
throughout, sliding steeply but continuously, with no fold. Shading marks bone
volume fractions below the linear-elastic validity floor, where the branch is
extrapolation. (B) The independent full-model probe with geometry free:
driving oestrogen down and restoring it recovers 100.3\,\% of baseline with no
hysteresis.}
\label{fig7}
\end{figure}

Freezing the slowly-evolving envelope geometry to isolate porosity dynamics,
and counting zeros of the bone-volume-fraction nullcline as oestrogen falls,
gives \textbf{one} fixed point throughout. It slides continuously from
$f_{bm}^* = 0.078$ at full oestrogen to 0.059 at $E_2 = 0.65$, passing steeply
but continuously through a critical value near $E_2 \approx 0.92$.
Mechanosensitivity and basal sclerostin secretion give the same answer, as
does scanning the strength of the osteocyte-density positive feedback across
1.5--5. There is no fold, no bistability and no hysteresis.

A second, independent route agrees. Driving oestrogen down and then restoring
it in the \textbf{full model with geometry free} --- a trajectory that stays
inside the linear-elastic domain end to end --- takes bone volume fraction
from 0.950 down to 0.836 and back to \textbf{0.953, that is 100.3\,\% of
baseline}, with no hysteresis.

The mechanism of monostability is instructive. Falling bone volume fraction
collapses the apparent modulus, strain rises steeply, mechanical dose and
modelling drift both increase, and the resulting recovery drive overwhelms the
two positive feedbacks we built in specifically to produce bistability --- loss
of remodelling surface as bone disappears \cite{martin1984}, and loss of
sensing capacity as osteocytes die. \textbf{The mechanostat rescues low bone
mass.}

One caution and one nuance. The deep branch of the continuation, below
$f_{bm} \approx 0.391$, lies outside the linear-elastic domain the mechanics
assumes \cite{bayraktar2004} and must be read as extrapolation; the
monostability conclusion depends on zero-counting and is unaffected. And the
recovery cycle leaves a permanent $+0.005$\,mm of periosteal expansion,
because modelling adds to the outer surface and nothing removes it. That is
structural memory --- physically real, since periosteal apposition is largely
irreversible \emph{in vivo} --- and it is emphatically not dynamical
bistability, since bone volume fraction recovers completely.

We note that a numerical artefact nearly produced the opposite conclusion. An
early full-system test showed bone volume fraction stuck at its floor for 40
years, which read as irreversibility. Inspection showed the ``collapsed''
state had a cortical thickness of 99\,mm: the modelling term was linear in
strain excess and therefore unbounded, and at pathological strain it demanded
1513\,mm/yr of periosteal apposition. Bounding it --- at a ceiling of
1.93\,$\mu$m/day, which the yield strain of cortical bone
\cite{bayraktar2004,morgan2001} and the documented maximum mineral apposition
rate independently agree on --- removed the artefact and made the
geometry-free probe above possible.

% ---------------------------------------------------------------------------
\section*{Discussion}

\subsection*{What the model buys}

The end-to-end chain lets three questions be asked that neither half of the
literature could answer alone, and each returns a number rather than a
direction: loading outweighs calcium 15.1-fold under stated conditions; a
systemic agent cannot produce a local asymmetry, structurally rather than
parametrically; and an anti-sclerostin antibody and mechanical loading compete
for a shared node to the tune of 13\,\% sub-additivity.

The result we would defend most strongly is the one that cost nothing. Frost's
set point, a phenomenological constant inferred from observation in 1987,
emerges here at 787\,\ue{} from channel gating kinetics and signalling Hill
functions, inside the band he proposed, without being fitted. Together with
the six held-out tennis-player indices and the held-out withdrawal loss, this
is evidence that the model predicts rather than accommodates.

\subsection*{Falsifying our own hypothesis}

P3 is the result we expected to confirm and did not. We report it in the main
text for two reasons.

First, it is more informative than confirmation would have been. Complex
physiological models exhibit bistability readily, and such findings are
fragile --- ours would have been, since a numerical artefact almost delivered
exactly that conclusion. What the falsification buys instead is a boundary:
the mechanostat's mechanical negative feedback is strong enough, at the
porosity scale, to rescue low bone mass, so clinical irreversibility
\textbf{cannot} originate in continuum remodelling dynamics. It must come from
mechanisms this model does not represent --- trabecular perforation, which is
loss of the template a continuous bone volume fraction cannot describe;
osteocyte death; or geometric degradation.

Second, that boundary is a testable prediction rather than a shrug. It says
that a compartment which can lose its template --- a trabecular one, where
struts perforate and cannot be rebuilt --- should show the bistability that a
cortical one does not. Our trabecular compartment is the natural place to
look, and we did not find bistability there either, because the compartment is
still described by a continuous bone volume fraction. Representing perforation
explicitly is the concrete next step.

We would also distinguish two things the literature often conflates. The
clinical asymmetry between rapid loss and slow regain is a \textbf{rate}
asymmetry, present in our monostable model, and it is not evidence of
dynamical irreversibility. Steepness is not bistability either: the transition
at $E_2 \approx 0.92$ is sharp enough to look like a jump, and the sensitivity
of disuse loss to the sclerostin Hill constant is likewise steep, but both are
high gain near a single fixed point.

\subsection*{Measurement choice is a result, not a detail}

One of P1's three clauses reverses sign under a different measure: the loading
effect is $+4.582\,\%$ in bone mineral content and $+1.778\,\%$ in areal
density. A second clause, the calcium--loading interaction, did the same until
the systemic calcium module was rebuilt, and is now positive on both framings
--- which makes the point more sharply, since the sign of a reported
interaction turned on a defect two modules away.

Neither reversal is a modelling error. The areal-density dilution is a
faithful reproduction of a known limitation \cite{haapasalo2000}: loading
translates the cortex outward, and dividing by projected width discards much
of the gain. The interaction reversal is a floor effect in the comparator,
since calcium deficiency harms the sedentary arm more than the loaded one.

The implication is practical. A trial powered on areal density will see
roughly two-fifths of the loading effect that a trial measuring bone mass
would see, and a trial reporting an interaction term may report the opposite
sign from a trial reporting achieved gains, from identical underlying biology.
We therefore state each clause of P1 with its measure and its loading
intensity attached, and we would encourage the same discipline in the
empirical literature.

\subsection*{Limitations}

The frequency response disagrees with experiment in sign \cite{hsieh2001},
traced to saturation in channel opening rate; the mechanosensitive channel
rate constants are placeholders taken from \cite{fu2025} and this is the first
quantity to re-check against better values. The logarithmic \emph{form} of the
frequency dependence is consistent with time-lapsed micro-computed tomography
estimates \cite{marques2023}; only the direction is not.

The specific surface function relating available remodelling surface to bone
volume fraction is a placeholder, and we can now say exactly how wrong it is.
Lerebours and colleagues \cite{lerebours2015} measured the relationship in human
cortical bone and report that at high bone volume fraction it behaves as
$(1-\mathrm{BV/TV})^{b}$ with $b = 0.5$, the square-root fit outperforming both
the classical fifth-order polynomial and a harmonic-mean form ($r^2 = 0.84$
against 0.82 and 0.77), and above $\mathrm{BV/TV} = 0.9$ their data converge on a
fractional degree rather than any integer one. \textbf{Our exponent is 3}, six
times the measured value, and our cortical compartment sits at a bone volume
fraction of 0.95 --- inside the range they measured. Specific surface here
therefore responds about six-fold too steeply to a change in porosity:
$d\ln S_v/d\ln f_{bm}$ is $-56$ at our cortical baseline where Lerebours implies
about $-9.5$. Baseline turnover is unaffected, since specific surface is
normalised at baseline, but every rate that passes through it is exaggerated.

\textbf{We report a negative result from trying to fix this, because it is more
informative than the fix would have been.} Setting the cortical exponent to the
measured 0.5 works mechanically --- the sensitivity becomes $-8.51$, close to
what the measurement implies --- but it drops the disuse target to
0.44\,\%/month, and the sclerostin sensitivity cannot recover it: scanning that
parameter to the bottom of its admissible range saturates at 0.99\,\%/month
without entering the 1.0--1.5\,\%/month band. A substantial part of the disuse
loss this model produces has therefore been carried by a specific-surface
response six times steeper than the real one, and with the measured exponent
something else must supply it.

\textbf{We should be precise about what that does and does not establish,
because our first statement of it was too strong.} Holding the turnover
constant fixed, as above, no setting of the sclerostin sensitivity reaches the
band. Releasing the turnover constant as well, a refit does reach
1.01\,\%/month --- but it puts adult turnover at 10.2\,\%/yr, outside the
5--10\,\%/yr band that is itself the upper half of what is conventionally
quoted \cite{parfitt2002}. The accurate claim is therefore not that the disuse
target is unreachable with the measured exponent. It is that every route we
have found pays for it with a different target.

We then built the mechanism that would have to supply it, and the result
changed twice --- which is why we report the sequence rather than only the
conclusion. The candidate is forced: the mechanical arm of sclerostin
production is already saturated in disuse, reaching 99.97\,\% of its absolute
ceiling after 180 days of simulated bed rest, so no route through
mechanosensing can supply further loss and the missing term must bypass
sclerostin entirely. Dying osteocytes releasing RANKL directly is such a route,
and we implemented it as a factor on the osteocyte deficit, exactly inert at
baseline.

Measured against the model as it then stood, the term worked and the objection
was its cost: it returned disuse loss to its band but made calcium deficiency
progressive, because the only thing that lowers osteocyte density here is
resorption itself. Measured again after the osteocyte turnover constant was
corrected (below), it does not work at all. Pushed to the top of its admissible
range it lifts disuse loss only from 0.44 to 0.50\,\%/month, nowhere near the
1.0--1.5 band, because the osteocyte deficit that drives it over 180 days of bed
rest is 2.2\,\% rather than the 20.7\,\% we had measured before --- and that
20.7\,\% was itself a product of the error. \textbf{The mechanism did not fail;
the evidence that it could succeed was an artefact.} We report both passes
because the first would have been a defensible published result and was wrong
for a reason no target would have caught.

Two independent checks agree with where this ended. Adding unloading-induced
osteocyte apoptosis, so that the deficit is generated mechanically rather than
by resorption, makes it disuse-specific as intended but changes nothing, and the
inductions it would need are far outside measurement: two weeks of hind-limb
unloading raises the proportion of apoptotic osteocytes by 66\,\%
\cite{basso2006}, where this route needs a fold two orders of magnitude larger.
And the experiment has been done. Blocking osteocyte apoptosis with a
bisphosphonate analogue that does not itself inhibit resorption prevents the
rise in osteocytic RANKL but does not prevent the resorption or the bone loss
caused by unloading \cite{plotkin2015}, from which the authors conclude that
RANKL from non-osteocytic sources carries the loss.

A third mechanism completes the search, and it is the one that came closest.
Immobilisation raises cortical \emph{porosity} --- that is resorption moving to
the intracortical surface, and it is measured in humans: the femoral cortex of
long-term immobilised individuals is significantly more porous, in a pattern
distinguishable from postmenopausal osteoporosis \cite{rolvien2020}. Written as
a bias on the resorption surface split driven by the mechanical dose deficit,
this route is disuse-specific by construction and by measurement --- the deficit
is 0.886 over 180 days of bed rest against 0.0016 in the low-calcium arm, so it
does not repeat the failure above. It also has the authority the apoptotic route
lacked: it lifts disuse loss from 0.44 to 1.04\,\%/month with the calcium
effect, its plateau and the romosozumab response all inside their bands.

\textbf{It fails on physiology instead.} Closing the gap requires a coefficient
at which 94.5\,\% of disuse resorption is intracortical and 5.0\,\%
endocortical, against a 50/45 baseline split --- that is nearly all resorption
moved to one surface, and human bed rest thins the endocortex as well as opening
porosity. A coefficient that only works where the state it produces is not
observed is a fitting parameter wearing a mechanism's name, and we decline to
report it as the resolution.

\textbf{Three routes have now failed in three different ways, and the pattern is
more informative than any of them.} The apoptotic arm was not scenario-specific,
because the deficit driving it is produced by resorption itself, and it cost the
calcium plateau. Making that deficit mechanical fixed the specificity and left
the route with no authority, since basal apoptosis is about 1\,\% of osteocyte
removal here. The surface bias has both specificity and authority and asks for a
surface split that is not observed. Specificity, authority, physiology: each
route fails a different one of the three, which is a stronger statement than
three failures of the same kind would be. On this evidence, reproducing disuse
loss at the measured specific-surface exponent needs a mechanism none of these
three supplies, and we think the most likely candidate is one this model does
not represent at all --- the loss of trabecular and endocortical \emph{template},
which is also where the falsification of P3 points.

One further consequence of the measured exponent deserves recording, because it
survived every pass: it makes disuse loss effectively irreversible. Ten years
of normal loading after 180 days of bed rest ends 5.7\,\% below baseline with no
net recovery, against a full recovery at the exponent we ship.

We therefore know what the exaggerated exponent was standing in for, and we know
that \textbf{the evidence it could be replaced was itself an artefact}. The
reason for shipping the value we do is no longer that the exchange is a bad
bargain; it is that the candidate route cannot supply the disuse loss at all
once the model is self-consistent. The measured 0.5 remains the sourced value and
we ship 3, and we record that plainly.

\textbf{A second and larger error came out of checking the first, and it was an
internal contradiction rather than a disagreement with a measurement.}
Osteocytes leave the tissue when the bone they occupy is resorbed, so the
fractional turnover of the osteocyte population must equal the fractional
turnover of bone. Until we checked, it did not. Bone was resorbed at
$1.93\times10^{-4}$ per day, the rate implied by our own adult turnover, while
osteocytes were removed at $9.9\times10^{-2}$ per day --- a factor of 514, and a
mean osteocyte residence of ten days rather than years. No external data are
needed to see this; the model disagreed with itself, and every target it had
passed for two years.

An independent estimate confirms which side was wrong: 42 billion osteocytes in
the adult skeleton with 9.1 million replenished daily \cite{buenzli2015} gives
$2.17\times10^{-4}$ per day, a residence of 12.6 years. We have corrected it,
and the correction is structural rather than a new value: the burial constant is
now \textbf{derived} from bone turnover at every parameter set, as the formation
constant already was, so the contradiction cannot recur. All results in this
paper are computed after that correction and a five-parameter recalibration.

The consequences run both ways. Against: the calcium effect fell to
$+0.663$\,\%, just outside its band, since the calcium arm had been amplified by
an osteocyte-density feedback that a ten-day residence made unphysiologically
fast. In favour: the post-withdrawal hold-out improved from finishing 0.75\,\%
below pre-treatment baseline to finishing 2.68\,\% above it, with no drug
parameter touched; ten-year recovery after bed rest became complete rather than
partial; and the reported PTH-pathway parameter left the upper bound it had
previously been pinned to. Baseline drift, the elastic domain and monostability
are unaffected.

We therefore report results computed with an exponent we know to be too steep,
and say so rather than quietly adopting the measured value and reporting a
broken disuse target. A further caution comes from the title of the same paper:
the relationship is \textbf{subject specific}, not intrinsic --- per-subject fits
reach $r^2$ of 0.87--0.98 against 0.84 pooled, correlating with sex and pore
density but not with age, height or weight. A single curve is an approximation
not only between compartments but between individuals.

Complete unloading beyond about 7.5 months leaves the model's valid domain, so
long-duration disuse is out of scope. The antibody is a step function rather
than a pharmacokinetic profile, which sharpens the withdrawal transient ---
the calibrated quantity is the downstream resorption overshoot, which is
unaffected in magnitude. Two PTH-pathway parameters are identifiable only
jointly. The cortical geometry is an idealised circular annulus.

Three limitations follow directly from the four targets reported above, and we
regard them as the most useful things a reader can take from this paper about
where it is weakest. The systemic calcium module was rebuilt during this
work, having failed in both directions at once: serum calcium swung 15\,\% with
dietary intake where physiology allows about 2\,\%, and the coefficients
carrying skeletal calcium into the serum pool were three orders of magnitude
below the flux they were normalised against, so the coupling was bidirectional
in structure but one-way in magnitude. Both are repaired --- the excursion is now
1.6\,\% and unloading suppresses parathyroid hormone by 21.3\,\% against a
reported 24\,\% \cite{spatz2012} --- but two consequences remain. The repair
cost the calcium effect size: V7 now sits at the low edge of its band rather
than mid-band, because the earlier value depended on the swing that has been
removed. And it created the V7/V16 conflict described above, in which no single
PTH$\rightarrow$RANKL gain satisfies both the calcium target and the
withdrawal-rebound target. Neither is hidden, and both are consequences of
making the module more physiological rather than less. The RANKL response to
sclerostin cannot reach its reported magnitude at any parameter value, because
the relevant half-saturation constant sits below baseline sclerostin. And the
model does not produce postmenopausal periosteal widening, so the geometry of
the ageing cortex is reproduced only in its thinning half. Finally, the
model has no representation of trabecular architecture as architecture, which
is precisely where the P3 falsification points.

\subsection*{Conclusion}

Joining a validated mechanotransduction front end to a validated bone-cell
back end produces a model that predicts blind observations, computes a set
point that was previously only observed, and settles the calcium-versus-loading
question quantitatively --- while overturning our own hypothesis about the
dynamical nature of osteoporosis, and showing that two of our three
quantitative claims depend on which measure is used to state them.

% ---------------------------------------------------------------------------
\section*{Materials and methods}

\subsection*{Model structure}

Eight modules; full equations in \nameref{S1_Text}. State vectors: 17 (single
compartment), 30 (two compartments; 13 local states per site plus four shared
systemic states). Integration with \texttt{ode15s}, relative tolerance
$1\times10^{-6}$, absolute tolerance $1\times10^{-9}$, non-negativity enforced
on all states.

The per-site biology is implemented once and called by both the single- and
two-compartment right-hand sides, so the two cannot diverge. Mechanosensitive
channel gating exists in exactly one place; the interface between fast
(sub-second) and slow (daily) timescales is a single scalar, the daily
integral of channel open probability. An earlier version represented the same
phenomena four times over and left that scalar with no downstream consumer, so
mechanical signalling never actually crossed the timescale boundary.

\subsection*{Loading specification}

Scenarios specify peak bending moment and peak axial force per loading bout,
with cycle count, frequency, inter-cycle rest, post-bout rest and days of
week. \textbf{Strain is never an input.} This is enforced by an automated
check that rejects any scenario carrying a strain-like field.

\subsection*{Calibration and validation}

Five free parameters fitted with \texttt{surrogateopt} (fixed seed) against five
targets; \texttt{k\_form} and \texttt{k\_ot} derived at each evaluation, from
bone-mass balance and from bone turnover respectively. Eight observations held
out: the six peripheral quantitative computed
tomography indices of the loaded humerus, the post-withdrawal areal density
change, and the emergent set point. Targets and tolerances are listed in
\nameref{S2_Table}; all parameters, with source and confidence grading, in
\nameref{S1_Table}. Hard-coded numerical constants in source are treated as
defects and are caught by an automated check.

\subsection*{Validity domain}

Every simulation reports whether peak strain remained below the yield strain
of cortical bone (7000\,\ue) \cite{bayraktar2004}, beyond which the linear
elasticity assumed throughout the mechanics stack fails. Results from runs
that leave this domain are not interpreted. This check is what identified the
long-duration disuse scope limit and the extrapolated deep branch of the
bifurcation diagram.

\subsection*{Bifurcation analysis}

Monostability was assessed two ways: by counting zeros of the
bone-volume-fraction nullcline with envelope geometry frozen, and
independently by a withdrawal-and-recovery trajectory in the full model with
geometry free. Fixed-point stability from Jacobian eigenvalues.

\subsection*{Software and reproducibility}

MATLAB R2026a. The model, the parameter and target tables, all seven
experiment scripts and a 91-assertion automated test suite are available at
https://github.com/myhsieh1002/bone-mechanostat and archived on Zenodo (concept
DOI 10.5281/zenodo.21592303; this version 10.5281/zenodo.21784609). Every figure in this
paper is regenerated by a single named script.

% ---------------------------------------------------------------------------
\section*{Supporting information}

\paragraph*{S1 Text.}
\label{S1_Text}
{\bf Model equations.} Full specification of modules M1--M8, state vector
layout, the trabecular compartment, the validity domain, and the derivation of
the bone-mass-balance constraint on the formation rate constant.

\paragraph*{S1 Table.}
\label{S1_Table}
{\bf Parameters.} Every parameter with value, unit, bounds, owning module,
literature source and confidence grading.

\paragraph*{S2 Table.}
\label{S2_Table}
{\bf Validation targets.} All targets with ranges, tolerances, associated
experiment and hold-out status.

\paragraph*{S1 Fig.}
\label{S1_Fig}
{\bf Identifiability.} Two-dimensional $\chi^2$ surface for the correlated
PTH-pathway pair and one-dimensional profiles for the three remaining fitted
parameters, of which two are individually identifiable and the antibody
clearance constant is not.

% ---------------------------------------------------------------------------
\section*{Data availability}

All model code, parameter tables, validation target definitions, experiment
scripts and the automated test suite are deposited at
https://github.com/myhsieh1002/bone-mechanostat and archived on Zenodo (concept
DOI 10.5281/zenodo.21592303; this version 10.5281/zenodo.21784609). No empirical data
were generated; all validation targets are drawn from published literature and
are listed with sources in S2 Table.

\section*{Author contributions}

PLOS collects contributions through the submission system using the CRediT
taxonomy; this statement mirrors what is entered there.
Hsiang-Lin Lee: Conceptualization; Investigation; Supervision; Validation; Writing -- review \& editing.
Tzu-Ling Wang: Data curation; Investigation; Resources; Writing -- review \& editing.
Ming-Yu Hsieh: Conceptualization; Data curation; Formal analysis;
Investigation; Methodology; Project administration; Software; Validation;
Visualization; Writing -- original draft; Writing -- review \& editing.

\section*{Acknowledgments}

The authors received no funding for this work.

\nolinenumbers

% ---------------------------------------------------------------------------
% To use BibTeX instead, comment out the thebibliography block below and
% uncomment these two lines (requires plos2015.bst from the PLOS template):
% \bibliographystyle{plos2015}
% \bibliography{references}

\begin{thebibliography}{10}

\bibitem{frost1987}
Frost HM.
\newblock Bone ``mass'' and the ``mechanostat'': a proposal.
\newblock Anat Rec. 1987;219(1):1--9. doi:10.1002/ar.1092190104.

\bibitem{fu2025}
Fu R, Wang C, Shahneela N, Ud Din R, Yang H.
\newblock A whole bone--lacunocanalicular network--osteocyte model examining
  bone adaptation to distinct loading parameters.
\newblock Int J Mech Sci. 2025;286:109931. doi:10.1016/j.ijmecsci.2025.109931.

\bibitem{lemaire2004}
Lemaire V, Tobin FL, Greller LD, Cho CR, Suva LJ.
\newblock Modeling the interactions between osteoblast and osteoclast
  activities in bone remodeling.
\newblock J Theor Biol. 2004;229(3):293--309. doi:10.1016/j.jtbi.2004.03.023.

\bibitem{pivonka2008}
Pivonka P, Zimak J, Smith DW, Gardiner BS, Dunstan CR, Sims NA, et~al.
\newblock Model structure and control of bone remodeling: a theoretical study.
\newblock Bone. 2008;43(2):249--263. doi:10.1016/j.bone.2008.03.025.

\bibitem{scheiner2013}
Scheiner S, Pivonka P, Hellmich C.
\newblock Coupling systems biology with multiscale mechanics, for computer
  simulations of bone remodeling.
\newblock Comput Methods Appl Mech Eng. 2013;254:181--196.

\bibitem{tai2015}
Tai V, Leung W, Grey A, Reid IR, Bolland MJ.
\newblock Calcium intake and bone mineral density: systematic review and
  meta-analysis.
\newblock BMJ. 2015;351:h4183. doi:10.1136/bmj.h4183.

\bibitem{haapasalo2000}
Haapasalo H, Kontulainen S, Siev\"{a}nen H, Kannus P, J\"{a}rvinen M, Vuori I.
\newblock Exercise-induced bone gain is due to enlargement in bone size without
  a change in volumetric bone density: a peripheral quantitative computed
  tomography study of the upper arms of male tennis players.
\newblock Bone. 2000;27(3):351--357. doi:10.1016/s8756-3282(00)00331-8.

\bibitem{peterson2010}
Peterson MC, Riggs MM.
\newblock A physiologically based mathematical model of integrated calcium
  homeostasis and bone remodeling.
\newblock Bone. 2010;46(1):49--63. doi:10.1016/j.bone.2009.08.053.

\bibitem{weinbaum1994}
Weinbaum S, Cowin SC, Zeng Y.
\newblock A model for the excitation of osteocytes by mechanical
  loading-induced bone fluid shear stresses.
\newblock J Biomech. 1994;27(3):339--360. doi:10.1016/0021-9290(94)90010-8.

\bibitem{srinivasan2002}
Srinivasan S, Weimer DA, Agans SC, Bain SD, Gross TS.
\newblock Low-magnitude mechanical loading becomes osteogenic when rest is
  inserted between each load cycle.
\newblock J Bone Miner Res. 2002;17(9):1613--1620.
  doi:10.1359/jbmr.2002.17.9.1613.

\bibitem{hsieh2001}
Hsieh YF, Turner CH.
\newblock Effects of loading frequency on mechanically induced bone formation.
\newblock J Bone Miner Res. 2001;16(5):918--924.
  doi:10.1359/jbmr.2001.16.5.918.

\bibitem{marques2023}
Marques FC, Boaretti D, Walle M, Scheuren AC, Schulte FA, M\"{u}ller R.
\newblock Mechanostat parameters estimated from time-lapsed in vivo
  micro-computed tomography data of mechanically driven bone adaptation are
  logarithmically dependent on loading frequency.
\newblock Front Bioeng Biotechnol. 2023;11:1140673.
  doi:10.3389/fbioe.2023.1140673.

\bibitem{schulte2026}
Schulte FA, Marques FC, Griesbach JK, Weigt C, von Salis-Soglio M, Lambers FM,
  et~al.
\newblock Combined physical and pharmacological anabolic osteoporosis therapies
  increase bone response and mechanoregulation in female mice.
\newblock Nat Commun. 2026;17:3759. doi:10.1038/s41467-026-70309-2.

\bibitem{cosman2016}
Cosman F, Crittenden DB, Adachi JD, Binkley N, Czerwinski E, Ferrari S, et~al.
\newblock Romosozumab treatment in postmenopausal women with osteoporosis.
\newblock N Engl J Med. 2016;375(16):1532--1543. doi:10.1056/NEJMoa1607948.

\bibitem{wijenayaka2011}
Wijenayaka AR, Kogawa M, Lim HP, Bonewald LF, Findlay DM, Atkins GJ.
\newblock Sclerostin stimulates osteocyte support of osteoclast activity by a
  RANKL-dependent pathway.
\newblock PLoS One. 2011;6(10):e25900. doi:10.1371/journal.pone.0025900.

\bibitem{currey1988}
Currey JD.
\newblock The effect of porosity and mineral content on the Young's modulus of
  elasticity of compact bone.
\newblock J Biomech. 1988;21(2):131--139. doi:10.1016/0021-9290(88)90006-1.

\bibitem{bayraktar2004}
Bayraktar HH, Morgan EF, Niebur GL, Morris GE, Wong EK, Keaveny TM.
\newblock Comparison of the elastic and yield properties of human femoral
  trabecular and cortical bone tissue.
\newblock J Biomech. 2004;37(1):27--35. doi:10.1016/s0021-9290(03)00257-4.

\bibitem{martin1984}
Martin RB.
\newblock Porosity and specific surface of bone.
\newblock Crit Rev Biomed Eng. 1984;10(3):179--222.

\bibitem{lerebours2015}
Lerebours C, Thomas CDL, Clement JG, Buenzli PR, Pivonka P.
\newblock The relationship between porosity and specific surface in human
  cortical bone is subject specific.
\newblock Bone. 2015;72:109--117. doi:10.1016/j.bone.2014.11.016.

\bibitem{morgan2001}
Morgan EF, Keaveny TM.
\newblock Dependence of yield strain of human trabecular bone on anatomic site.
\newblock J Biomech. 2001;34(5):569--577.

\bibitem{yang2017}
Yang H, Embry RE, Main RP.
\newblock Effects of loading duration and short rest insertion on cancellous
  and cortical bone adaptation in the mouse tibia.
\newblock PLoS One. 2017;12(1):e0169601.

\bibitem{li2019}
Li X, Han L, Nookaew I, Mannen E, Silva MJ, Almeida M, et~al.
\newblock Stimulation of Piezo1 by mechanical signals promotes bone anabolism.
\newblock eLife. 2019;8:e49631.

\bibitem{spatz2012}
Spatz JM, Fields EE, Yu EW, Divieti Pajevic P, Bouxsein ML, Sibonga JD, et~al.
\newblock Serum sclerostin increases in healthy adult men during bed rest.
\newblock J Clin Endocrinol Metab. 2012;97(9):E1736--40.

\bibitem{parfitt2002}
Parfitt AM.
\newblock Misconceptions (2): turnover is always higher in cancellous than in
  cortical bone.
\newblock Bone. 2002;30(6):807--809.

\bibitem{basso2006}
Basso N, Heersche JNM.
\newblock Effects of hind limb unloading and reloading on nitric oxide synthase
  expression and apoptosis of osteocytes and chondrocytes.
\newblock Bone. 2006;39(4):807--814. doi:10.1016/j.bone.2006.04.014

\bibitem{plotkin2015}
Plotkin LI, Gortazar AR, Davis HM, Condon KW, Gabilondo H, Maycas M, et al.
\newblock Inhibition of osteocyte apoptosis prevents the increase in osteocytic
  receptor activator of nuclear factor $\kappa$B ligand (RANKL) but does not
  stop bone resorption or the loss of bone induced by unloading.
\newblock J Biol Chem. 2015;290(31):18934--18942. doi:10.1074/jbc.M115.642090

\bibitem{rolvien2020}
Rolvien T, Milovanovic P, Schmidt FN, von Kroge S, W\"olfel EM, Krause M, et al.
\newblock Long-term immobilization in elderly females causes a specific pattern
  of cortical bone and osteocyte deterioration different from postmenopausal
  osteoporosis.
\newblock J Bone Miner Res. 2020;35(7):1343--1351. doi:10.1002/jbmr.3970

\bibitem{buenzli2015}
Buenzli PR, Sims NA.
\newblock Quantifying the osteocyte network in the human skeleton.
\newblock Bone. 2015;75:144--150. doi:10.1016/j.bone.2015.02.016

\bibitem{mcclung2018}
McClung MR, Brown JP, Diez-Perez A, Resch H, Caminis J, Meisner P, et~al.
\newblock Effects of 24 months of treatment with romosozumab followed by 12
  months of denosumab or placebo in postmenopausal women with low bone mineral
  density: a randomized, double-blind, phase 2, parallel group study.
\newblock J Bone Miner Res. 2018;33(8):1397--1406.

\end{thebibliography}

\end{document}
